\chapter{Conclusión}\label{ch:conclusion}

El presente trabajo abordó la falta de asistencia automática, gratuita y adaptada al contexto argentino para evaluar la veracidad de una publicación en el momento en que se la lee. Se desarrolló un prototipo que combina un clasificador de lenguaje natural, las señales públicas de la cuenta autora y el contraste con evidencia externa, entregado como extensión de Google Chrome para Twitter/X con un panel web complementario. Las conclusiones recorren los cuatro objetivos específicos del Capítulo~1 y cierran con las limitaciones del prototipo y las líneas de trabajo futuro.

En relación con el primer objetivo, el protocolo de evaluación del clasificador se fijó antes de obtener cualquier resultado, para que la elección de los datos y de las métricas no dependiera de los valores observados. La estrategia de datos de la Sección~\ref{sec:datos} combina la transferencia desde corpus en inglés con la adaptación al español, y la Sección~\ref{sec:evaluacion-clasificador} establece el valor F1 macro como métrica principal frente a una línea base clásica, dos modelos monolingües y un modelo de lenguaje grande sin ajuste, que mide el costo de oportunidad de entrenar un modelo propio. \Martin[inline]{Agregar el F1 macro sobre el conjunto de prueba real, la diferencia con la línea base y con el modelo de lenguaje grande, y si se cumple RNF-05.}

El segundo objetivo, referido a las señales de la cuenta autora, se vio modificado por la investigación con usuarios. Las entrevistas de la Sección~\ref{sec:entrevistas} mostraron que esas señales tienen un valor informativo alto, porque son lo primero que consulta quien verifica, pero una validez probatoria baja, porque la antigüedad, los seguidores y la verificación pueden adquirirse. Por ese motivo el módulo se conserva como información exhibida al usuario y pesa menos que el clasificador y que el contraste externo en el puntaje (Sección~\ref{sec:arquitectura}). \Martin[inline]{Describir la versión mínima implementada y su prueba. Si no se implementa, declarar el objetivo como cumplido de forma parcial.}

El tercer objetivo concentra el aporte diferencial del sistema. El módulo de contraste semántico recupera evidencia de tres clases de fuente ordenadas por jerarquía: las fuentes oficiales de un índice local actualizado por ingesta programada, los medios del padrón de socios activos de ADEPA, adoptado en lugar de una selección propia para no exponer el sistema a la sospecha de sesgo, y las verificaciones previas de organizaciones profesionales. Cada fuente se clasifica según su postura y se enlaza con el fragmento que la sostiene, y las pruebas automatizadas de la Sección~\ref{sec:pruebas-automatizadas} verifican que ningún veredicto se emita sin una fuente enlazable. \Martin[inline]{Agregar la proporción de análisis con al menos tres fuentes relevantes, meta que fija el objetivo.}

El cuarto objetivo, la integración en una extensión y un panel web, se verificó con las pruebas funcionales de la Sección~\ref{sec:pruebas-funcionales}, una por caso de uso, que cubren quince de los dieciséis requerimientos funcionales. Las baterías automatizadas suman 61 pruebas en el servicio, con una cobertura del 88~\% de las sentencias, y 16 en la extensión, y aseguran que la falla de una etapa se comunique como análisis parcial. La prueba de usabilidad de la Sección~\ref{sec:usabilidad} aporta la evidencia sobre la interpretación del indicador por personas sin conocimiento técnico. \Martin[inline]{Agregar los casos de prueba aprobados, la latencia en el percentil 95 y el SUS medio. El objetivo fija menos de 5~s y RNF-02, 8~s en el percentil 95: unificar el umbral.}

El prototipo presenta limitaciones que conviene declarar. La principal concierne a la naturaleza del puntaje: la verdad de una afirmación no admite grados, y el valor que muestra la extensión no expresa una probabilidad de verdad sino el grado en que la evidencia recuperada apoya la afirmación extraída. Una afirmación compuesta por una parte verdadera y otra falsa es falsa por conjunción, pero el ensamblador promedia la postura de las fuentes y puede ubicarla en un nivel intermedio. A ello se suma que la extracción conserva una sola afirmación verificable por publicación, de modo que las restantes no se evalúan. Asimismo, la jerarquía otorga el mayor peso a la fuente oficial, con lo cual un dato oficial erróneo se traslada al veredicto sin que el sistema pueda advertirlo, y un razonamiento que parte de una premisa falsa puede derivar cualquier conclusión, según el principio \textit{ex falso quodlibet}. El diseño acota estos riesgos sin resolverlos: exhibe cada fuente con su postura, atribuye el nivel más severo a la fuente que contradice y no al sistema, y emite el estado \textit{sin contraste externo} cuando no encuentra evidencia. Por último, la encuesta de la Sección~\ref{sec:encuestas} se apoya en un muestreo no probabilístico concentrado entre los 18 y los 25 años, por lo que sus hallazgos no se extrapolan al conjunto de usuarios de Twitter/X.

Las líneas de trabajo futuro se organizan en tres ejes, que corresponden a formas de engaño que el prototipo todavía no detecta. El primero es el razonamiento sobre la afirmación. En lugar de promediar puntajes, la afirmación puede descomponerse en afirmaciones atómicas, cada una con una sola pieza de información, como en FActScore \parencite{MinEtAl2023}, o en subpreguntas explícitas e implícitas cuyas respuestas determinan su veracidad, como en ClaimDecomp \parencite{ChenEtAl2022}, y el veredicto puede derivarse de un programa de razonamiento que resuelve cada subtarea con un verificador especializado, como en ProgramFC \parencite{PanEtAl2023}. Este enfoque permitiría evaluar todas las afirmaciones de una publicación, marcar como falsa una conjunción con un solo componente contradicho y limitar un dato oficial erróneo al predicado que lo contiene, a cambio de multiplicar la latencia y el costo por publicación.

El segundo eje es el contexto omitido: una publicación puede engañar con datos verdaderos si calla la información que los contradice o si recircula como novedad un contenido antiguo, y tanto la periodista entrevistada como las respuestas abiertas de la encuesta pidieron incorporar la fecha de origen y el contexto faltante. El tercer eje es el contenido visual: un gráfico con el eje recortado convierte una diferencia menor en un salto aparente sin que ningún dato sea falso, y detectarlo exige leer imágenes, un alcance que el Capítulo~1 excluyó por su costo.

En síntesis, el trabajo muestra que es posible acercar al ciudadano, en el momento de la lectura, la evidencia necesaria para juzgar una publicación. El propósito del sistema no es que el usuario crea en el veredicto, sino que pueda prescindir de él, y por eso cada resultado se acompaña del camino hacia las fuentes que lo sostienen.
